\section{Hipótesis de Trabajo}
 
\subsection{Hipótesis de Primer Grado}
 
Los antipatrones arquitectónicos definidos en la taxonomía propuesta se encuentran presentes en aplicaciones de código abierto basadas en Inteligencia Artificial Generativa y RAG.

Esta hipótesis es de carácter descriptivo y será contrastada mediante la aplicación del instrumento de análisis estático sobre la muestra seleccionada. Su comprobación permitirá establecer si los patrones definidos conceptualmente pueden ser identificados en implementaciones reales del dominio estudiado.
 
\subsection{Hipótesis de Segundo Grado}
 
La presencia de antipatrones arquitectónicos en aplicaciones basadas en IA Generativa y RAG está asociada con un mayor deterioro de las métricas de mantenibilidad y calidad estructural del software, particularmente en términos de acoplamiento entre objetos (CBO), complejidad ciclomática e índice de mantenibilidad, en comparación con componentes en los que dichos antipatrones no se encuentren presentes.

Esta hipótesis establece la relación entre la variable independiente, representada por la presencia y frecuencia de antipatrones arquitectónicos, y las variables dependientes asociadas con las métricas de calidad estructural y mantenibilidad. Su contrastación se realizará mediante el análisis estadístico de los datos obtenidos de los repositorios seleccionados.
 
\subsection{Hipótesis de Tercer Grado}
 
No se plantea una hipótesis de tercer grado para el presente estudio, debido a que los objetivos de investigación se concentran en la caracterización de los antipatrones, su detección automatizada y el análisis de su relación con métricas tradicionales de acoplamiento, complejidad y mantenibilidad. Incorporar variables adicionales, como el framework empleado, la antigüedad del repositorio o el número de colaboradores, ampliaría el alcance del estudio más allá de las relaciones definidas en el problema y los objetivos de investigación.